\documentclass[oneside]{scrbook}

% \usepackage[utf8]{inputenc}
\usepackage[T1]{fontenc}
\usepackage{version}
% \usepackage[latin1]{inputenc}
\usepackage[italian]{babel}
\usepackage{algpseudocode}
\usepackage{amsmath}
\usepackage{booktabs}
\usepackage{color}
\usepackage{amsfonts}
\usepackage{url}
\usepackage{algorithm}
\usepackage{chngcntr}
\usepackage[answerdelayed]{exercise}
\usepackage{graphicx}
\usepackage{caption}
\usepackage{listings}
\usepackage{tikz}
\usepackage{pgf}
\usepackage{colortbl}
\usepackage{listingsutf8}
\usepackage[pdftitle = {Dispense di esercizi}, pdfauthor = {Gianmarco Tedeschi e Alberto Marchesi}, pdfsubject = {Fondamenti di Informatica ($082746$ - A.A.~$2025-2026$)}, pdfcreator = {Gianmarco Tedeschi e Alberto Marchesi}, pdfproducer = {-}]{hyperref}
% \usepackage[automark,nouppercase]{scrpage2}
\usepackage{mdframed}
\usepackage{scrlayer-scrpage} % Necessary package for page styles

%\pagestyle{scrheadings}

%\usepackage{titlesec}
%\newcommand{\sectionbreak}{\clearpage}
\counterwithin{Exercise}{chapter}
\counterwithin{Answer}{chapter}

\algrenewcommand\algorithmicthen{\textbf{Allora}}
\algrenewcommand\algorithmicif{\textbf{Se}}
\algrenewcommand\algorithmicelse{\textbf{Altrimenti}}
\algrenewcommand\algorithmicend{\textbf{Chiudi}}
\algrenewcommand\algorithmicwhile{\textbf{Finch\'e}}
\algrenewcommand\algorithmicdo{\textbf{Esegui}}
\algrenewcommand{\algorithmiccomment}[1]{\hskip3em\texttt{//#1}}

\setlength{\parindent}{0in}
\setlength{\parskip}{10pt}

\newmdenv[%
        outerlinewidth=2,%
        roundcorner=10pt,%
        leftmargin=0,%
        rightmargin=0,%
        outerlinecolor=black,
        innertopmargin=\topskip,%
        splittopskip=\topskip,%
]{algo}


\ohead{\headmark}
\ihead{}
\chead{}
\ifoot{A.A.~$2025-2026$}
\cfoot{Fondamenti di Informatica (082746) - G. Tedeschi}
\ofoot{Pag.~\pagemark}
\setheadsepline{.4pt}
\setfootsepline{.4pt}

\definecolor{lbcolor}{rgb}{0.9,0.9,0.9}

\lstdefinestyle{cstyle}{%
  inputencoding=utf8/latin1,
  breaklines=true,
  breakatwhitespace=false,
  basicstyle=\small\ttfamily,
  language=c,
  frame=single,
  showstringspaces=false
}

\lstdefinestyle{cstylesmall}{%
  inputencoding=utf8/latin1,
  breaklines=true,
  breakatwhitespace=false,
  basicstyle=\scriptsize\ttfamily,
  language=c,
  frame=single,
  showstringspaces=false
}

\lstdefinestyle{console}{%
  inputencoding=utf8/latin1,
  breaklines=true,
  breakatwhitespace=false,
  basicstyle=\scriptsize\ttfamily,
  showstringspaces=false,
  backgroundcolor=\color{lbcolor}
}


\lstnewenvironment{clang}{\lstset{style=cstyle}}{}
\lstnewenvironment{clangsmall}{\lstset{style=cstylesmall}}{}

\addtocounter{chapter}{6}

\begin{document}
\chapter{Funzioni}

Questa dispensa propone esercizi sulla scrittura di algoritmi, in linguaggio C, utili alla comprensione della definizione e dell'utilizzo delle \textbf{funzioni}.

Le funzioni sono blocchi di codice riutilizzabili che eseguono un compito specifico. Permettono di organizzare il programma in parti più piccole e gestibili (principio di \emph{modularità}), migliorando la leggibilità e riducendo la duplicazione del codice.

Un programma C è composto da una o più funzioni; la funzione \texttt{main()} è obbligatoria ed è il punto di ingresso del programma.

\section{Dichiarazione e Definizione}

Per usare una funzione, di solito sono necessari due passaggi: la \textbf{dichiarazione} (prototipo) e la \textbf{definizione} (implementazione).

\paragraph{Dichiarazione (Prototipo)}

La dichiarazione, o prototipo, informa il compilatore sull'esistenza di una funzione, specificando il suo nome, il tipo di dato che restituisce (tipo di ritorno) e i tipi di dati che accetta come input (parametri).

La sintassi è:
\begin{lstlisting}[style=cstyle]
tipo_ritorno nome_funzione(tipo_param1, tipo_param2, ...);
\end{lstlisting}

I prototipi si inseriscono tipicamente all'inizio del file, dopo gli \texttt{\#include}.

\paragraph{Definizione}

La definizione è il blocco di codice che costituisce la funzione vera e propria. Contiene le istruzioni che verranno eseguite quando la funzione è chiamata.

La sintassi è:
\begin{lstlisting}[style=cstyle]
tipo_ritorno nome_funzione(tipo_param1 nome_param1, tipo_param2 nome_param2, ...) {
    
// Corpo della funzione
    // ... istruzioni ...
    
    return valore_restituito; // Solo per funzioni non void
}
\end{lstlisting}
\begin{itemize}
    \item \textbf{\texttt{tipo\_ritorno}}: Il tipo di dato del valore che la funzione restituisce al chiamante. Se la funzione non restituisce nulla, si usa \texttt{void}.
    \item \textbf{\texttt{nome\_funzione}}: Il nome della funzione.
    \item Parametri: Le variabili (es. \texttt{nome\_param1}) che ricevono i valori passati alla funzione.
    \item \textbf{\texttt{return}}: L'istruzione che termina la funzione e (se non \texttt{void}) restituisce un valore.
\end{itemize}

\section{Scope delle Variabili}

Lo \emph{scope} (visibilità) di una variabile definisce dove nel programma essa può essere utilizzata. Ogni variabile esiste solo all'interno del suo scope.

\begin{itemize}
    \item \textbf{Variabili Locali}: Dichiarate all'interno di una funzione (inclusi i parametri). Sono visibili e utilizzabili \textbf{solo} all'interno di quella funzione. Vengono create quando la funzione inizia e distrutte (perdendo il loro valore) quando la funzione termina. Le variabili locali di una funzione (es. \texttt{main}) non sono visibili da un'altra funzione (es. \texttt{swap\_errato}).
    \item \textbf{Variabili Globali}: Dichiarate all'esterno di qualsiasi funzione (di solito in cima al file). Sono visibili e modificabili da \textbf{qualsiasi} funzione nel file (dal punto di dichiarazione in poi). Il loro uso è generalmente sconsigliato perché rendono il programma difficile da comprendere e manutenere, creando dipendenze nascoste tra le funzioni.
\end{itemize}


\section{Passaggio dei Parametri}

\paragraph{Passaggio per Valore}

"Passare per valore" significa che, quando una funzione viene chiamata, il \textbf{valore} dell'argomento nel chiamante (es. \texttt{main}) viene \textbf{copiato} in una \emph{nuova} variabile: il parametro formale della funzione chiamata.

La funzione lavora quindi esclusivamente su questa \textbf{copia locale}. Qualsiasi modifica fatta a questa copia \emph{all'interno} della funzione \textbf{non} ha alcun effetto sulla variabile originale nel chiamante.

Ad esempio, un'operazione di scambio tra due varabili effettuata con il passaggio per valore non avrà effetto al di fuori della funzione:

\begin{lstlisting}[style=cstyle]
// 'a' e 'b' sono copie locali, esistono solo qui dentro.
void swap_errato(int a, int b) {
    int temp = a;
    a = b;
    b = temp;
    // Qui 'a' (la copia) e 'b' (la copia) sono state scambiate.
    // Le variabili originali 'x' e 'y' del main non vengono toccate.
}

void main() {
    int x = 10;
    int y = 20;
    
    // Il valore di x (10) viene copiato in 'a'
    // Il valore di y (20) viene copiato in 'b'
    swap_errato(x, y);
    
    // Al termine di swap_errato, 'a' e 'b' vengono distrutte.
    // 'x' e 'y' non sono mai state modificate.
    // Il risultato della printf rimane ancora "10, 20"
    printf("%d, %d", x, y); 
}
\end{lstlisting}

\paragraph{Passaggio per Riferimento}

Cosa fare se vogliamo che una funzione modifichi le variabili del chiamante? 
Possiamo utilizzare \textbf{l'indirizzo di memoria} della variabile come parametro formale.

Il meccanismo è sempre un passaggio "per valore", ma il valore che passiamo non è il valore della variabile (es. 10), ma il valore del suo indirizzo (es. 0x7FFF... ).

\begin{itemize}
    \item Nel chiamante (\texttt{main}): Usiamo l'operatore \texttt{\&} per ottenere l'indirizzo della variabile (es. \texttt{\&x}).
    \item Nella funzione: Il parametro formale deve essere di tipo \textbf{puntatore} (dichiarato con \texttt{*}) per poter memorizzare un indirizzo.
    \item Nella funzione: Per accedere o modificare il valore \emph{contenuto} a quell'indirizzo, usiamo l'operatore di "dereferenziazione" (\texttt{*}).
\end{itemize}

Possiamo ora effettuare uno scambio tra variabili che verrà propagato anche nel main:

\begin{lstlisting}[style=cstyle]
void swap(int *a, int *b) {
    int temp;
    
    // *a significa "vai all'indirizzo contenuto in a"
    // e prendi il valore (10).
    temp = *a; 

    // *b significa "vai all'indirizzo 'b' (a 'y')"
    // e prendi il valore (20).
    // *a = ... significa "vai all'indirizzo 'a' (a 'x') e
    // scrivi il nuovo valore (20).
    *a = *b;       
    
    // *b = ... significa "vai all'indirizzo 'b' (a 'y') e 
    // scrivi il valore di temp (10).
    *b = temp;       
}

void main() {
    int x = 10;
    int y = 20;
    
    // Passo l'indirizzo di x 
    // Passo l'indirizzo di y 
    swap(&x, &y); 
    
    // La funzione ha modificato direttamente i valori
    // alle locazioni di memoria di x e y.
    // Ora x vale 20 e y vale 10.
    printf("%d, %d", x, y);
}
\end{lstlisting}

\section{Funzioni Ricorsive}
Le funzioni ricorsive sono funzioni che si richiamano direttamente o indirettamente. La ricorsione è utile per risolvere problemi che possono essere suddivisi in sottoproblemi simili. Per risolvere un problema attraverso la programmazione ricorsiva sono necessari alcuni elementi:
\begin{itemize}
    \item Caso base: caso elementare del problema che può essere
    risolto immediatamente.
    \item Passo ricorsivo: chiamata ricorsiva per risolvere uno o più problemi più semplici.
    \item Costruzione della soluzione: costruzione della soluzione
    sulla base del risultato delle chiamate ricorsive
\end{itemize}
Durante la costruzione di funzioni ricorsive è fondamentale aggiungere una condizione di terminazione che ci assicurerà che prima o poi la catena di chiamate a funzioni termini. 

\textbf{Esempio:} Calcolo del Fattoriale. Il fattoriale di un numero \( n \) (\( n! \)) può essere calcolato ricorsivamente come:
\[
n! = 
\begin{cases} 
1 & \text{se } n = 0, \\
n \cdot (n-1)! & \text{se } n > 0.
\end{cases}
\]
\begin{lstlisting}[style=cstyle]
int fattoriale(int n) {
    if (n == 0)        // Caso base
        return 1;
    else               // Caso ricorsivo
        return n * fattoriale(n - 1);
}
\end{lstlisting}



\section{Esercizi}

\begin{Exercise}[label={ex:operazioni}]
    Scrivere un programma che legga una sequenza di numeri interi positivi, terminata dal valore 0.
    Il programma deve calcolare, il numero totale di valori inseriti (escluso lo 0), la somma di tutti i numeri,
    il massimo e il minimo valore inserito e la media dei valori inseriti.

    Scrivere le funzioni:
    \begin{itemize}
        \item leggiNumero(): che legge e restituisce un numero naturale
        \item massimo(): restituisce il maggiore tra due numeri
        \item minimo(): restituisce il minore tra due numeri
        \item calcolaMedia(): restituisce la media dei numeri letti
    \end{itemize}
\end{Exercise}

\begin{Answer}[ref={ex:operazioni}]
  	\lstinputlisting[style=cstylesmall]{src/operazioni.c}
\end{Answer}
\begin{Exercise}[label={ex:amici}]
    Scrivere un programma che prende in input due numeri naturali $n1$ e $n2$ e verifica se sono amici
    Due numeri si dicono amici quando la somma dei divisori propri di $n1$ è uguale a $n2$ e viceversa.
    \begin{itemize}
        \item leggiNumero(): che legge e restituisce un numero naturale
        \item numeroAmico(): che riceve in ingresso due numeri naturali e restituisce vero
        \item sommaDivisori(): che riceve in ingresso un numero naturale e restituisce la somma dei suoi divisori (eccetto se stesso)
    \end{itemize}
\end{Exercise}

\begin{Answer}[ref={ex:amici}]
  	\lstinputlisting[style=cstylesmall]{src/amici.c}
\end{Answer}
\begin{Exercise}[label={ex:perfetti}]
    Scrivere un programma che chiede all’utente di inserire un numero naturale. Il programma verifica che il numero sia perfetto.
    Un numero si dice perfetto quando è uguale all somma di tutti i suoi divisori (eccetto se stesso).

    Scrivere le funzioni:
    \begin{itemize}
        \item leggiNumero(): che legge e restituisce un numero naturale
        \item numeroPerfetto(): che riceve in ingresso un numero naturale e restituisce vero
        \item sommaDivisori(): che riceve in ingresso un numero naturale e restituisce la somma dei suoi divisori (eccetto se stesso)
    \end{itemize}
\end{Exercise}

\begin{Answer}[ref={ex:perfetti}]
  	\lstinputlisting[style=cstylesmall]{src/perfetti.c}
\end{Answer}
\begin{Exercise}[label={ex:sequenza}]
	Scrivere una funzione che riceve un array di dimensione $N$ e un intero $k$ e calcola la lunghezza della sequenza piu' lunga di interi consecutivi distanti tra loro esattamente $k$.

	\textbf{Esempio}: $k=3$, sequenza: 1 2 5 8 1 $\rightarrow$ 2
\end{Exercise}

\begin{Answer}[ref={ex:sequenza}]
  	\lstinputlisting[style=cstylesmall]{src/sequenza.c}
\end{Answer}
\begin{Exercise}[label={ex:ricerca}]
    Si realizzi un programma che legga da input un numero non noto a priori di interi (al più 10) terminati da 0.
    Tale sequenza potrà evantualmente contenere numeri ripetuti.
    A tal fine, si realizzi una funziona \texttt{leggi(int v[], int dim)} che riceva in ingresso un vettore (vuoto) e la sua dimensione fisica,
    e restituisca la dimensione logica.
    
    Si chieda poi l'inserimento di un valore da cercare nell'array. 
    Si realizzi una funzione \texttt{search(int v[], int dim, int k)} che stampi a video tutti gli indici relativi alle posizioni in cui il numero k sia presente nell’array. 
\end{Exercise}

\begin{Answer}[ref={ex:ricerca}]
  	\lstinputlisting[style=cstylesmall]{src/ricerca.c}
\end{Answer}
\begin{Exercise}[label={ex:ripetizioni}]
    Creare un programma che legga da input una sequenza di interi, di lunghezza non nota a priori (al più 10), e terminata da 0. A tal scopo, si
    realizzi una funzione che riceva come parametri di ingresso un vettore e la sua dimensione fisica, e restituisca la dimensione logica del
    vettore. Tale funzione si deve fare carico della fase di lettura e riempimento dell’array. La sequenza può contenere numeri ripetuti.

    Si realizzi una funzione che, ricevuti in ingresso il primo vettore con la sua dimensione logica, ed un secondo vettore con la sua dimensione
    fisica, memorizzi nel secondo vettore tutti gli elementi del primo, ma senza ripetizioni. La funzione restituisca la dimensione logica del
    secondo vettore.

    Si realizzi un main che invochi le funzioni, e che stampi a video l’elenco degli elementi non ripetuti
\end{Exercise}

\begin{Answer}[ref={ex:ripetizioni}]
  	\lstinputlisting[style=cstylesmall]{src/ripetizioni.c}
\end{Answer}
\begin{Exercise}[label={ex:rimuovi}]
    Scrivere un programma che legge una stringa senza spazi e un carattere e stampa la stringa inserita senza il carattere inserito.
    Scrivere le funzioni:
    \begin{itemize}
        \item leggiCarattere(): che legge e restituisce un carattere
        \item leggiStringa(): che legge e restituisce una stringa
        \item rimuoviCarattere(): che prende in input una stringa e un carattere e restituisce la stringa senza il carattere inserito
    \end{itemize}
\end{Exercise}

\begin{Answer}[ref={ex:rimuovi}]
  	\lstinputlisting[style=cstylesmall]{src/rimuovi.c}
\end{Answer}
\begin{Exercise}[label={ex:shift}]
    Scrivere un programma che legge una matrice 3x4 ed effettua uno shift per righe verso il basso e stampa il risultato. L’ultima riga prende la posizione della prima riga.
    Scrivere le funzioni:   
    \begin{itemize}
        \item leggiArray(): che legge e restituisce un array di 4 elementi
        \item stampaArray(): che stampa un array di 4 elementi
        \item copiaArray(): che copia un array di 4 elementi in un altro array di 4 elementi
        \item leggiMatrice(): che legge e restituisce una matrice 3x4
        \item stampaMatrice(): che stampa una matrice 3x4
        \item shiftMatrice(): che effettua lo shift della matrice verso il basso
    \end{itemize}
\end{Exercise}

\begin{Answer}[ref={ex:shift}]
  	\lstinputlisting[style=cstylesmall]{src/shift.c}
\end{Answer}

%Riferimento
\begin{Exercise}[label={ex:operazioni_rif}]
    Si scriva una funzione che riceva due varibili (base ed esponente) per valore, e due variabili di output (potenza e radice)
    passate per riferimento. La funzione deve calcolare la potenza (base elevata ad esponente) e la radice (radice ennesima di base) e restituirle
    tramite le variabili di output.
\end{Exercise}

\begin{Answer}[ref={ex:operazioni_rif}]
  	\lstinputlisting[style=cstylesmall]{src/operazioni_rif.c}
\end{Answer}
\begin{Exercise}[label={ex:punto}]
    Si scriva un programma che gestisca un punto nel piano cartesiano.
    Si definisca una struttura \texttt{Punto} che rappresenti un punto nel piano.
    Si implementino le seguenti funzionalità:
    \begin{itemize}
        \item Una funzione \texttt{trasli} che riceve un punto e due valori \texttt{dx} e \texttt{dy} e trasla il punto di tali valori.
        \item Una funzione \texttt{baricentro} che riceve un insieme di punti e restituisce il punto baricentro.
    \end{itemize}
\end{Exercise}

\begin{Answer}[ref={ex:punto}]
  	\lstinputlisting[style=cstylesmall]{src/punto.c}
\end{Answer}

% Ricorsione
\begin{Exercise}[label={ex:ricorsione_1}]
    Si scrivano le versioni ricorsiva ed iterativa di una funzione \texttt{double f(double a, int n)}
    che calcoli il seguente valore:
    $$
        \sum_{i=1}^{n} a - \frac{i}{a}
    $$
\end{Exercise}

\begin{Answer}[ref={ex:ricorsione_1}]
  	\lstinputlisting[style=cstylesmall]{src/ricorsione_1.c}
\end{Answer}
\begin{Exercise}[label={ex:ricorsione_2}]
    Si scriva un programma che inverta le cifre di un numero intero $N$ usando una funzione apposita.
    Si realizzi sia una versione ricorsiva che una versione iterativa.

    \textbf{Esempio:} se l'input è 1234, l'output deve essere 4321.
\end{Exercise}

\begin{Answer}[ref={ex:ricorsione_2}]
  	\lstinputlisting[style=cstylesmall]{src/ricorsione_2.c}
\end{Answer}
\begin{Exercise}[label={ex:inv_ric}]
    Si scriva una funzione ricorsiva che inverta una stringa data in ingresso.
\end{Exercise}

\begin{Answer}[ref={ex:inv_ric}]
  	\lstinputlisting[style=cstylesmall]{src/inv_ric.c}
\end{Answer}
\begin{Exercise}[label={ex:hanoi}]
    La \emph{torre di Hanoi} è un famoso rompicapo matematico composto da tre aste "\texttt{A}", "\texttt{B}" e "\texttt{C}" e un numero \texttt{N} di dischi di grandezza decrescente.
    Il rompicapo inizia con i dischi impilati in ordine crescente di dimensione sulla prima asta (\texttt{A}), il disco più piccolo in cima.
    L'obiettivo del gioco è spostare l'intera pila di dischi dalla prima asta (\texttt{A}) alla terza asta (\texttt{C}), rispettando le seguenti regole:
    \begin{itemize}
        \item Si può spostare solo un disco alla volta.
        \item Non si può posizionare un disco più grande sopra uno più piccolo.
    \end{itemize}
    Si scriva una funzione che risolva il problema in modo ricorsivo per un numero \texttt{N} di dischi dato in input,
    stampando ad ogni mossa il disco che viene spostato, l'asta di partenza e l'asta di destinazione.

    \textbf{Esempio:} Per \texttt{N = 3}, l'output dovrebbe essere:
    \begin{verbatim}
    Sposta il disco 1 da A a C
    Sposta il disco 2 da A a B
    Sposta il disco 1 da C a B
    Sposta il disco 3 da A a C
    Sposta il disco 1 da B a A
    Sposta il disco 2 da B a C
    Sposta il disco 1 da A a C
    \end{verbatim}
\end{Exercise}

\begin{Answer}[ref={ex:hanoi}]
  	\lstinputlisting[style=cstylesmall]{src/hanoi.c}
\end{Answer}

% Complicati
\begin{Exercise}[label={ex:triathlon}]
    Si scriva un programma per la gestione degli atleti che partecipano a una gara di triathlon.
    Ogni atleta è caratterizzato da un nome (stringa di massimo 50 caratteri), da un identificativo e un tempo totale di gara (in ore, minuti e secondi).
    Si implementino le seguenti funzionalità:
    \begin{itemize}
        \item Definizione delle strutture dati necessarie per rappresentare un atleta e una lista di atleti.
        \item Scrivere una funzione \texttt{somma\_tempo} che riceve in input due tempi di gara e restituisce la somma dei due tempi.
        \item Scrivere una funzione \texttt{verifica\_tempo\_minore} che ricevuti due tempi di gara restituisce 1 se la prima variabile è minore o uguale alla seconda, 0 altrimenti.
        \item Scrivere una funzione \texttt{calcola\_tempo\_totale} che riceve in input un atleta e calcola il tempo totale di tutta la competizione.
        \item Scrivere una funzione \texttt{stampa\_miglior\_tempo} che riceve in input un'array di atleti e la sua dimensione e stampa a schermo l'id e il cognome dell'atleta con il tempo tempo totale migliore.
    \end{itemize}
\end{Exercise}

\begin{Answer}[ref={ex:triathlon}]
  	\lstinputlisting[style=cstylesmall]{src/triathlon.c}
\end{Answer}

\clearpage
\section*{Soluzioni}
\shipoutAnswer

\end{document}